\documentclass[11pt, a4paper]{article}

% --- Packages ---
\usepackage[margin=1in]{geometry}
\usepackage{amsmath, amssymb, amsfonts}
\usepackage{mathpazo} % Professional font
\usepackage{microtype} % Improves typography

\begin{document}

\noindent \textbf{Question:} Find the minimum number of terms required to evaluate the Maclaurin series of $e^x$ for $x = 1$ such that the absolute error is $\le 0.0005$, using Digital Signal Processing (DSP) and Matrix Operator theory.

\vspace{0.6cm}

\noindent \textbf{Solution:} 

\vspace{0.3cm}

\noindent Let $\mathbf{S}$ be the infinite-dimensional downshift matrix, representing a discrete unit delay operator ($z^{-1}$). When multiplied by a vector, it shifts all elements down by one position:
\begin{equation}
    \mathbf{S} = \begin{bmatrix} 0 & 0 & 0 & \dots \\ 1 & 0 & 0 & \dots \\ 0 & 1 & 0 & \dots \\ \vdots & \vdots & \ddots & \ddots \end{bmatrix}, \quad \mathbf{S} \begin{bmatrix} a \\ b \\ c \\ \vdots \end{bmatrix} = \begin{bmatrix} 0 \\ a \\ b \\ \vdots \end{bmatrix}
\end{equation}

\noindent The ideal infinite impulse response of this system is generated by applying the matrix exponential to an impulse vector $\mathbf{\delta} = [1, 0, 0, \dots]^T$. Expanding the Taylor series of $e^{\mathbf{S}}$ directly generates the Maclaurin coefficients:
\begin{equation}
    \mathbf{h}_{\text{ideal}} = e^{\mathbf{S}} \mathbf{\delta} = \left( \mathbf{I} + \mathbf{S} + \frac{1}{2!} \mathbf{S}^2 + \frac{1}{3!} \mathbf{S}^3 + \dots \right) \begin{bmatrix} 1 \\ 0 \\ 0 \\ \vdots \end{bmatrix} = \begin{bmatrix} 1 \\ 1 \\ 1/2! \\ 1/3! \\ \vdots \end{bmatrix}
\end{equation}

\noindent Truncating this infinite system to a finite filter length $N$ yields the error tail matrix, $\mathbf{E}_N$:
\begin{equation}
    \mathbf{E}_N = \sum_{k=N}^{\infty} \frac{1}{k!} \mathbf{S}^k
\end{equation}

\noindent To measure the maximum size of this error, we apply the induced $l_1$-norm ($\|\cdot\|_1$), defined as the maximum absolute column sum of a matrix. Because every column of $\mathbf{S}$ contains exactly one $1$ and zeros elsewhere, its column sum is strictly $1$, hence $\|\mathbf{S}\|_1 = 1$. By the submultiplicative property of norms ($\|\mathbf{A}\mathbf{B}\| \le \|\mathbf{A}\|\|\mathbf{B}\|$), we establish $\|\mathbf{S}^k\|_1 \le \|\mathbf{S}\|_1^k = 1^k = 1$. This allows us to bound the error tail:
\begin{equation}
    \|\mathbf{E}_N\|_1 \le \sum_{k=N}^{\infty} \frac{\|\mathbf{S}\|_1^k}{k!} = \sum_{k=N}^{\infty} \frac{1}{k!}
\end{equation}

\noindent Expanding the summation and factoring out the first neglected term, $\frac{1}{N!}$, yields:
\begin{align}
    \|\mathbf{E}_N\|_1 &\le \frac{1}{N!} + \frac{1}{(N+1)!} + \frac{1}{(N+2)!} + \dots \\
    \|\mathbf{E}_N\|_1 &\le \frac{1}{N!} \left( 1 + \frac{1}{N+1} + \frac{1}{(N+1)(N+2)} + \dots \right)
\end{align}

\noindent To construct a solvable algebraic bound, we replace the increasing denominators with the worst-case constant decay ratio, $N$. Because $N+1 > N$, it is mathematically guaranteed that $\frac{1}{N+1} < \frac{1}{N}$. This substitution converts the complex fractional tail into a standard infinite geometric series:
\begin{equation}
    \|\mathbf{E}_N\|_1 \le \frac{1}{N!} \left( 1 + \frac{1}{N} + \frac{1}{N^2} + \dots \right) = \frac{1}{N!} \sum_{m=0}^{\infty} \left(\frac{1}{N}\right)^m
\end{equation}
\begin{equation}
    \|\mathbf{E}_N\|_1 \le \frac{1}{N!} \left( \frac{1}{1 - 1/N} \right) = \frac{1}{N!} \left( \frac{N}{N-1} \right)
\end{equation}

\noindent We enforce the target error tolerance ($\epsilon \le 0.0005$) on our derived bound:
\begin{equation}
    \frac{1}{N!} \left( \frac{N}{N-1} \right) \le 0.0005
\end{equation}

\noindent Evaluating the boundary conditions for the filter length $N$:
\begin{align}
    N &= 6: \quad \frac{1}{720} \left( \frac{6}{5} \right) = \frac{1}{600} \approx 0.00166 \quad (\text{Tolerance not met}) \\
    N &= 7: \quad \frac{1}{5040} \left( \frac{7}{6} \right) = \frac{7}{30240} \approx 0.00023 \quad (\text{Tolerance strictly met})
\end{align}

\vspace{0.3cm}

\noindent The minimum required number of terms is exactly $\mathbf{N = 7}$.

\end{document}
